
\subsection{Tessera's lung gain is matched by averaging at the scale of the annotation}
Table~\ref{tab:inputs} records observation and coordinate-frame coverage; reference-class counts and retained PCA variance describe target and molecular-input complexity (Supplementary Fig.~S1). Across the five NSCLC patients, mean ARI was \NsclcPca{} for PCA50 and \NsclcFull{} for Tessera; Tessera was higher in \LungVsPcaPositive{} of 5 patients (mean difference \NsclcDelta{}, exact \(p=\LungVsPcaP{}\); \panelref{fig:comparison}{a--b}, Supplementary Table~S16). The untuned six-neighbor mean reached \NsclcMean{}, within \NsclcVsMean{} of Tessera. When the neighborhood scale was selected on the other four patients, the selected control was a neighbor mean over 100 or 200 neighbors in every fold and reached mean ARI \LungNestedControlMean{}. It exceeded Tessera in all five patients (mean \LungNestedDelta{}, range \LungNestedMin{} to \LungNestedMax{}; exact \(p=\LungNestedP{}\); \panelref{fig:comparison}{b}, Supplementary Table~S17). Agreement rose steadily with neighborhood size up to the 200-neighbor scale at which the regions were constructed (Supplementary Fig.~S8). Complete ARI and NMI seed summaries are reported in Supplementary Tables S1--S4.

Whole-section sensitivity analyses of \WholeCellsMin{} to \WholeCellsMax{} cells reproduced this ordering with the fixed cohort-modal controls. Tessera exceeded PCA50 in \LungWholeVsPcaPositive{} of 5 lung sections (mean \LungWholeVsPcaMean{}), and the 200-neighbor mean exceeded Tessera in all five (mean \LungWholeVsControlMean{}; descriptive \(p=\LungWholeVsControlP{}\); \panelref{fig:comparison}{c}). This ordering persists without the 6,000-cell budget, but the same patients informed the control choice.







\subsection{Removing gates does not reduce lung agreement; contrast carries the gain}
If the gates detected tissue boundaries, removing them should reduce agreement and they should separate annotation seams. Neither prediction held consistently. Removing the gates increased lung agreement in four of five patients; Lung6 decreased by less than 0.001 ARI (mean change \NoGateMean{}; \panelref{fig:ablations}{a}). Gate seam AUROC ranged from \GateAurocMin{} to \GateAurocMax{} (median \GateAurocMedian{}), and a plain PCA-feature difference on the same edges separated seams better in \GateBeaten{} of \GateObjects{} objects (median AUROC \GateExprAurocMedian{}; \panelref{fig:ablations}{b}). Most last-layer lung gates were near zero (median \LungGateMedianMin{}--\LungGateMedianMax{}; \LungGateLowMin{}--\LungGateLowMax{}\% below 0.05), with strongly overlapping within- and across-region distributions (\panelref{fig:ablations}{c}; Supplementary Table S6). Nevertheless, their normalized relative weights were not uniform: patient-level seed means of median node total-variation distance from uniform ranged from \LungWeightTvdMin{} to \LungWeightTvdMax{} (Supplementary Table S8). The evidence for dispensable gating is the ablation, not an algebraic equivalence between small gates and an unweighted mean.

Message passing without the contrastive term added nothing: the no-contrast encoder matched the self-only encoder within \NoContrastSelfMaxGap{} ARI in every lung patient, and training on permuted coordinates without the contrastive term matched it within \PermNoContrastMaxGap{}. The lung gain therefore comes from the contrastive objective, which pulls spatial neighbors together in the embedding (\panelref{fig:ablations}{a}). Measured against the self-only encoder, which is the appropriate reference for the effect of geometry, full Tessera was better in four of five lung patients and TVA, lower in Lung6, and worse in \BelowSelfOnly{} of 13 objects, including all four BCC patients (\panelref{fig:ablations}{d}). The earlier comparison with permuted coordinates overstated the benefit of geometry, because permutation with the contrastive term retained makes random cells positives and degrades the embedding.



The lung display fields show consolidation of the supplied regions in Lung5, Lung9, Lung12 and Lung13 (\panelref{fig:lung}{a,c--e}). Lung6's nearly single-region display field highlights why field identity can dominate that patient's target (\panelref{fig:lung}{b}).



\subsection{Cell-type targets and clustering resolution change the comparison}
BCC references are molecular cell types. PCA50 reached mean ARI \BccPca{} and Tessera \BccFull{}; Tessera was lower in all four patients (mean \BccDelta{}; exact \(p=\BccVsPcaP{}\)), although the difference in Patient A was within seed variation (Supplementary Table~S16). The tuned control selected on the other BCC patients was a BANKSY-style representation that keeps each cell's own features and adds a lightly weighted 100-neighbor mean (\(\lambda=0.2\)). It exceeded Tessera in all four patients on the bounded fields (mean \BccNestedDelta{}) and on whole sections (mean \BccWholeVsControlMean{}; \panelref{fig:comparison}{b--c}). Increasing the neighborhood scale did not help BCC (Supplementary Fig.~S8), and the BCC display fields show intermingled categories for all four patients (Supplementary Fig.~S11).

The lung data allow the same comparison without changing dataset. Re-clustering the saved lung embeddings at the cell-type \(K\) and scoring them against supplied cell types, mean ARI was \CellTypePca{} for PCA50, \CellTypeTessera{} for Tessera and \CellTypeMean{} for the neighbor mean; Tessera was below PCA50 in \CellTypeTesseraLower{} of 5 patients (\panelref{fig:comparison}{d}). The same embeddings that gain on regions therefore lose on cell types, linking the direction of the effect to the target within one platform. With \(K\) set by classes of at least 1\% prevalence, but ARI still scored on all valid labels, Tessera exceeded PCA50 in \LungCommonKBetter{} of 5 lung patients (mean \LungCommonKDelta{}) and \BccCommonKBetter{} of 4 BCC patients (mean \BccCommonKDelta{}; Supplementary Table S7). The BCC ordering is thus sensitive to clustering resolution, not uniformly unfavorable across \(K\). At the primary \(K\), breast Visium, Open-ST and MIBI favored PCA50, and colorectal TVA favored Tessera (\panelref{fig:comparison}{a}). Colorectal assignments and measured protein intensities provide context for that comparison (Supplementary Fig.~S3), without independent validation labels. The breast ROI shows Tessera consolidating mixed spots into larger domains than the dominant-category reference (Supplementary Fig.~S12). Restricting evaluation to purer spots raised every arm's agreement but did not change their order (Supplementary Fig.~S13). The seam-score scatter against composition entropy and the supplied tumor/DCIS fractions retain the spot-level context (Supplementary Fig.~S13); they are inspection readouts, not independently validated boundaries. Open-ST and MIBI comparisons retain identical coordinates across arms (Supplementary Figs.~S5 and S6). Measured lung and BCC RNA markers provide tissue context (Supplementary Fig.~S7) without supplying independent validation labels.









\subsection{Patient transfer trades cell-type accuracy for region accuracy}
In leave-one-patient-out transfer, Tessera raised region balanced accuracy from \TransferRegionPca{} with PCA50 to \TransferRegionFull{}, but lowered cell-type balanced accuracy from \TransferCellPca{} to \TransferCellFull{}, in all five folds (\panelref{fig:transfer}{a--b}). Spatial neighbor means showed the same trade-off more strongly: the 200-neighbor mean reached \TransferRegionMeanTwoHundred{} for regions and \TransferCellMeanTwoHundred{} for cell types. Accuracy followed the same pattern (\panelref{fig:transfer}{c--d}). Tessera therefore sits between expression-only and smoothed representations, and the transfer setting reproduces the target dependence seen within objects. No supplied target category was absent from the sources in these folds. Supplementary Table S5 reports the complete patient-held-out accuracy and balanced-accuracy values.





\subsection{Tessera partitions vary between seeds}
For PCA50, the three seeds varied only the K-means initialization; the minimum seed-pair ARI was \StabilityPcaMin{}. Tessera's seed-pair ARI ranged from \StabilityTesseraMin{} to \StabilityTesseraMax{} (Supplementary Fig.~S14), so its partitions changed substantially between seeds even when mean agreement with the reference was similar. Agreement gains over PCA50 were confined to objects in which neighboring cells usually share a label, the lung regions and TVA (Supplementary Fig.~S14). Lung6 is the exception: its display field is almost entirely one region, and the region task there reduces to separating two fields. Execution summaries and saved loss curves document the observation budgets and fitting records (Supplementary Fig.~S4).

\subsection{An independent pathologist-labelled cohort reproduces the ordering}
Supplementary Table S9 reports annotation coverage and class composition. In the independent cohort (\HerPatients{} patients, \HerSpotsMin{} to \HerSpotsMax{} spots, \HerClassesMin{} to \HerClassesMax{} annotation classes), Tessera was below PCA50 in \HerAllTessPcaNeg{} patients (mean \HerAllTessPcaMean{}) and below the tuned control in \HerAllTessNestedNeg{} (mean \HerAllTessNestedMean{}; exact \(p=\HerAllTessNestedP{}\); \panelref{fig:external}{a}, Supplementary Table~S18). The untuned same-graph mean was higher than Tessera in \HerAllTessUSixNeg{} of \HerPatients{} patients (Tessera minus mean, \HerAllTessUSixMean{}; exact \(p=\HerAllTessUSixP{}\), bootstrap interval \HerAllTessUSixLo{} to \HerAllTessUSixHi{}) and higher than the tuned control in \HerAllUSixNestedPos{}.

Scored only on the tumor boundary, the same-graph mean exceeded the tuned control in all \HerPatients{} patients (mean \HerBndUSixNestedMean{}, exact \(p=\HerBndUSixNestedP{}\), interval \HerBndUSixNestedLo{} to \HerBndUSixNestedHi{}), whereas Tessera and the control split \HerBndTessNestedPos{} to \HerBndTessNestedNeg{} (mean \HerBndTessNestedMean{}) and Tessera was below the same-graph mean in \HerBndTessUSixNeg{} patients (\panelref{fig:external}{b,c}). The highest Tessera scores were in patients B, G and H, where the same-graph mean was also high; in patients C, E and F every method stayed below 0.14. The direction held for named compartments (same-graph mean above the tuned control in \HerCompUSixNestedPos{} of \HerPatients{}; Tessera below the mean in \HerCompTessUSixNeg{}) and for both variants of the boundary definition (Supplementary Table S10).

The maps show the two patients with the highest and the two with the lowest same-graph agreement (Supplementary Fig.~S9): in B and H the methods recover the main tumor masses with differing border roughness, and in C and E none recovers the annotated region closely.

The learned gate had tumor-boundary seam AUROC near chance (AUROC \HerBndSeamGate{}, 95\% bootstrap interval \HerBndSeamGateLo{} to \HerBndSeamGateHi{}; above 0.5 in \HerBndGateAboveChance{} of \HerPatients{} patients), and the plain feature difference was barely better (\HerBndSeamRaw{}). The exploratory protocol prediction that the gate would exceed the feature difference in at most three patients \HerHTwoVerdict{} (it did in \HerBndGateOverRaw{}), although both mean AUROCs were close to 0.5. Smoothing the feature difference over each neighborhood did separate it (AUROC \HerBndSeamSmooth{}, interval \HerBndSeamSmoothLo{} to \HerBndSeamSmoothHi{}) and exceeded the gate in \HerBndSeamSmoothGatePos{} of \HerPatients{} patients (mean \HerBndSeamSmoothGateMean{}, interval \HerBndSeamSmoothGateLo{} to \HerBndSeamSmoothGateHi{}; \panelref{fig:external}{d}); weighting the smoothed difference added \HerBndSeamSmoothPlainMean{}, a small edge-level gain that did not carry through to partitions (Supplementary Table S11). Repeated fits of the frozen encoder changed per-patient scores by up to \RepMaxRange{} ARI (median \RepMedianRange{}) over the \RepPairs{} patient-target pairs and moved Tessera to the other side of the tuned control in \RepSignChangesControl{} of them; no comparison with the same-graph mean changed sign (\panelref{fig:external}{e}, Supplementary Table S13), and the same-graph mean and tuned controls are deterministic.

Reweighting neighbors by expression similarity, which we tried to test whether an informed local rule could help, did not. The weighted mean was above the unweighted mean in \HerBndEUSixPos{} patients and below it in \HerBndEUSixNeg{} (mean \HerBndEUSixMean{}); permuting its weights across edges gave no consistent partition advantage (above the shuffled-weight version in \HerBndEShufPos{} of \HerPatients{}), whereas the weighted mean exceeded the coordinate-permuted version in \HerBndECoordPos{} of \HerPatients{} (mean \HerBndECoordMean{}; \panelref{fig:external}{f}). The 13 bounded objects agreed (weighted mean above the unweighted mean in \DevAllThirteenEMeanPos{} and below it in \DevAllThirteenEMeanNeg{}; mean \DevAllThirteenEMeanMean{}). These contrasts support the measured-neighborhood contribution; they do not establish equivalence between weighted and unweighted means.





\subsection{Tessera's advantage depends on the tissue and the annotation}
Three non-tumor tissues with anatomical labels behaved differently (Supplementary Table S12). Tessera exceeded PCA50 in all three, by \BrPcaGainMin{} to \BrPcaGainMax{}. It exceeded the tuned control in the human cortex section (\BrCortexTessNested{}) but not in mouse visual cortex (\BrStarmapTessNested{}) or hypothalamus (\BrMerfishTessNested{}), and only the cortex section exceeded the best of the 30 configurations chosen with the test labels (\BrCortexTessOracle{}, against \BrStarmapTessOracle{} and \BrMerfishTessOracle{}). The cortex section was used when the encoder was developed and first benchmarked, so it does not show that the learned component generalizes. Gate seam AUROC was below chance in all three (\BrGateMin{} to \BrGateMax{}).

Across all \LandN{} objects (\panelref{fig:landscape}{a--b}), Tessera exceeded the same-graph six-neighbor mean in \LandTessSamePos{} and the best configuration chosen with the labels in \LandTessOraclePos{}, the cortex section and colorectal TVA. Its advantage over the same-graph mean was positive in \LandComputedSamePos{} of \LandComputedN{} computed lung regions, \LandAnatSamePos{} of \LandAnatN{} anatomical tissues, \LandCellSamePos{} of \LandCellN{} cell-type labels and \LandPathSamePos{} of \LandPathN{} pathologist-labelled patients. The neighborhood size at which a uniform mean peaked depended on the annotation: \LandComputedKMin{} to \LandComputedKMax{} neighbors for the computed lung regions, \LandAnatKMin{} to \LandAnatKMax{} for anatomical tissues, \LandCellKMin{} to \LandCellKMax{} for cell-type labels, and \LandPathKMin{} to \LandPathKMax{} for the pathologist-labelled patients, whose best sizes split between local and large (\panelref{fig:landscape}{c}), so a scale selected on other patients is a compromise. Gate seam AUROC had median \LandGateMedian{} across the objects and exceeded 0.5 in \LandGateAbove{} of \LandN{}; the plain feature difference exceeded the gate in \LandRawOverGate{} (\panelref{fig:landscape}{d}).



\subsection{Untouched colorectal pathology supports the control, with weak gate discrimination}
In the QC-filtered colorectal pilot, the same-graph mean exceeded Tessera in \CrcTessMeanNeg{} of \CrcPatients{} patients (Tessera minus mean \CrcTessMeanDelta{} ARI; descriptive patient-bootstrap interval \CrcTessMeanLow{} to \CrcTessMeanHigh{}; Supplementary Table S14). The original target including the QC class gave the same ordering in \CrcRawTessMeanNeg{} of \CrcPatients{} patients (Supplementary Table S15). Removing gates changed agreement by \CrcNoGateDelta{} ARI on average (interval \CrcNoGateLow{} to \CrcNoGateHigh{}), which does not establish equivalence or a useful gating gain. The 18 complete model fits took \CrcFitSeconds{} CPU seconds; the neighborhood control required no model fitting.

The gate's mean seam AUROC was \CrcGate{} (interval \CrcGateLow{} to \CrcGateHigh{}), against \CrcRaw{} for the raw feature difference. It passed the pilot's above-chance bootstrap rule but did not demonstrate superiority over that simple score (gate minus raw \CrcGateRawDelta{}; interval \CrcGateRawLow{} to \CrcGateRawHigh{}). Thus, gates are not universally devoid of seam information. The new patients support the narrower finding that this information did not yield an advantage over a trivial edge score or a same-graph neighborhood mean. These are broad pathologist compartment seams, including mixed regions, not a validated binary tumor-boundary detector. The small, resource-selected pilot and its schema-triggered sensitivity provide additional external evidence rather than a decisive population-level confirmation.


